\documentclass[12pt,a4paper]{article}
\usepackage{geometry}
\geometry{left=2.5cm,right=2.5cm,top=2.0cm,bottom=2.5cm}
\usepackage[english]{babel}
\usepackage{amsmath,amsthm}
\usepackage{amsfonts}
\usepackage[longend,ruled,linesnumbered]{algorithm2e}
\usepackage{fancyhdr}
\usepackage{ctex}
\usepackage{array}
\usepackage{listings}
\usepackage{color}
\usepackage{graphicx}

\begin{document}


\title{
{《优化理论与算法》第8次作业
}
}
\date{}

% \author{
% 姓名：\underline{你的名字}~~~~~~
% 学号：\underline{你的学号}~~~~~~}

\maketitle
\begin{itemize}
	\item 作业截止于{\bf 周日（5/3/2026）}，在Canvas系统中提交PDF或照片文件
	\item 作业题目的tex文件可以在Canvas中下载
	\item 鼓励相互讨论，但不要抄袭.
\end{itemize}
%%%%%%%%%%%%%%%%%%%%%%%%%%%%%%%%%%%%%%%%%%%%%%%%%%%%%%%%%%%%%%%%%
%%%%%%%%%%%%%%%%%%%%%%%%%%%%%%%%%%%%%%%%%%%%%%%%%%%%%%%%%%%%%%%%%

\section{阅读与积累}
\paragraph{1.1}
\begin{itemize}
	\item 复习 Boyd and Vandenberghe，也就是 Convex Optimization 教材，章节 3.2-3.4 与 4.1-4.2相关
\end{itemize}

\section{习题}

\paragraph{2.1}
 判断下列函数是否是凸函数、凹函数、拟凸函数？
\begin{itemize}
  \item[(a)] $f(x) = e^x - 1$, 定义域为 $\mathbb{R}$
  \item[(b)] $f(x_1, x_2) = x_1 / x_2$, 定义域为 $\mathbb{R}_{++}^2$
  \item[(c)] $f(x_1, x_2) = x_1^2 / x_2$, 定义域为 $\mathbb{R} \times \mathbb{R}_{++}$
\end{itemize}



\paragraph{2.2} 复合规则。证明下列函数是凸函数：
\begin{itemize}
  \item[(a)] $f(x) = -\log\left(-\log\left(\sum_{i=1}^m e^{a_i^T x + b_i}\right)\right)$, \quad $\operatorname{dom} f = \left\{ x \mid \sum_{i=1}^m e^{a_i^T x + b_i} < 1 \right\}$
  \item[(b)] $f(x,u,v) = -\sqrt{uv - x^T x}, \quad \operatorname{dom} f = \{(x,u,v) \mid uv > x^T x, u > 0, v > 0\}$
\end{itemize}

\paragraph{2.3} 正式证明共轭函数为凸函数，即使原函数是非凸函数。

\paragraph{2.4} 求给定负熵函数的共轭：
\[
  f(x) = \sum_{i=1}^n x_i \log(x_i / 1^T x), \quad \operatorname{dom} f = \mathbb{R}_{++}^n
\]
证明其共轭函数为
\[
  f^*(y) =
  \begin{cases}
    0 & \text{若 } \sum_{i=1}^n e^{y_i} \leq 1 \\
    +\infty & \text{其他情况}
  \end{cases}
\]




\paragraph{2.5} 《Convex Optimization》(Boyd and Vandenberghe)教材课后题4.1

\paragraph{2.6} 《Convex Optimization》(Boyd and Vandenberghe)教材课后题4.3

\paragraph{2.7} 《Convex Optimization》(Boyd and Vandenberghe)教材课后题4.26 (a) （提示：利用上图形式，考虑目标函数形式为：$\min \; \sum_{i=1}^m t_i$）


\paragraph{2.8} 《Convex Optimization》(Boyd and Vandenberghe)教材课后题4.28 (a) （提示：利用上图形式，考虑目标函数形式为：$\min \; t$）

\end{document}